\begin{datatype}\label{propform}
    %! infixr 0
    Define $\propform$ inductively as follows.
    \begin{enumerate}
        \item $\propbot \in \propform$.
        \item $\propvar{n} \in \propform$ for
            $n \in \{\emptyset\}$.
        \item $(p \propto q) \in \propform$ for $p \in \propform$ and $q \in \propform$.
    \end{enumerate}
\end{datatype}
\begin{proposition}\label{propform_bot_test}
    $\propbot \in \propform$.
\end{proposition}
\begin{proof}
    Follows by \auto{propform_propbot_intro}.
\end{proof}

\begin{proposition}\label{propform_var_test}
    If $\emptyset \in \{\emptyset\}$, then
        $\propvar{\emptyset} \in \propform$.
\end{proposition}
\begin{proof}
    Follows by \auto{propform_propvar_intro}.
\end{proof}

\begin{proposition}\label{propform_imp_test}
    $(\propbot \propto \propbot) \in \propform$.
\end{proposition}
\begin{proof}
    Follows by \auto{propform_propbot_intro,propform_propto_intro}.
\end{proof}

\begin{proposition}\label{propform_distinct_test}
    $\propbot \neq (\propbot \propto \propbot)$.
\end{proposition}
\begin{proof}
    Follows by \auto{propform_propbot_propto_distinct}.
\end{proof}

\begin{proposition}\label{propform_injective_test}
    If $\propvar{x} = \propvar{y}$, then $x = y$.
\end{proposition}
\begin{proof}
    Follows by \auto{propform_propvar_injective}.
\end{proof}
